\chapter*[Introdução]{Introdução}\label{chap:intro}
\addcontentsline{toc}{chapter}{INTRODUÇÃO}

A robótica móvel consolidou-se como um campo de convergência entre computação, eletrônica e controle, no qual sistemas físicos passaram a depender de camadas de software cada vez mais elaboradas para percepção, decisão e comunicação. Em plataformas de pequeno porte, é comum que o robô não opere de forma totalmente autônoma, mas sim por teleoperação, com um operador humano enviando comandos de movimento à distância enquanto o sistema registra continuamente informações sobre o próprio estado -- posição, deslocamento, consumo de energia e conectividade. Esse registro contínuo, conhecido como telemetria, transforma o robô também em uma fonte de dados reais, que, uma vez armazenados e analisados, permitem compreender e validar o comportamento da solução construída.

O tratamento adequado desses dados, porém, não é imediato. À medida que um protótipo evolui de peças isoladas para uma plataforma integrada, cresce a complexidade do software que a sustenta, e com ela a necessidade de organizá-lo segundo princípios que favoreçam a modularidade, a reutilização e a testabilidade -- preocupações centrais da Programação Orientada a Objetos e da Engenharia de Software \cite{sommerville2018}. Do mesmo modo, dados de telemetria só se tornam informação útil quando persistidos em uma estrutura de armazenamento coerente com sua natureza temporal e submetidos a técnicas estatísticas que revelem padrões de comportamento. Esses três aspectos -- estruturação orientada a objetos, armazenamento de dados e análise estatística -- articulam-se sob práticas de gerenciamento de projetos que ordenam o desenvolvimento ao longo do tempo.

É nesse cenário que se insere o robô Atlas, uma plataforma robótica móvel desenvolvida no curso de Engenharia de Computação da SETREM. O robô já opera como um sistema completo, controlado por um aplicativo móvel e capaz de publicar sua telemetria em um serviço de nuvem que a persiste em banco de dados; entretanto, os dados coletados ainda não eram submetidos a nenhuma etapa de análise estatística. Este trabalho parte desse sistema em funcionamento para consolidá-lo como uma solução computacional orientada a objetos e, sobre a base de dados já em operação, estruturar a coleta, o armazenamento e a análise da telemetria, aplicando de forma integrada os quatro componentes curriculares do Projeto Interdisciplinar de Extensão IV.

O presente documento está organizado em três capítulos. O primeiro apresenta a contextualização do estudo, incluindo tema, objetivos, problema, justificativa, metodologia, cronograma e orçamento, sob a orientação da disciplina de Gerenciamento de Projetos. O segundo aborda a fundamentação teórica que sustenta as escolhas de projeto orientado a objetos, de modelagem de banco de dados e de análise estatística, encerrando com a descrição da arquitetura da solução. O terceiro detalha os resultados obtidos até a presente etapa, contemplando o desenvolvimento da solução computacional, a coleta e o armazenamento dos dados e as análises realizadas.
